% Handwritten numerical expectations; no empirical results.
\documentclass[10pt]{article}
\usepackage[a4paper,margin=17mm]{geometry}
\usepackage[T1]{fontenc}
\usepackage{lmodern,booktabs,tabularx,array,amsmath,xcolor,fancyhdr,draftwatermark,float}
\usepackage[font=small,labelfont=bf]{caption}
\newcolumntype{L}[1]{>{\raggedright\arraybackslash}p{#1}}
\newcolumntype{Y}{>{\raggedright\arraybackslash}X}
\newcommand{\rare}{\textsc{rare}}
\newcommand{\rng}[1]{{\footnotesize[#1]}}
\newcommand{\stack}[2]{\begin{tabular}[c]{@{}c@{}}#1\\[-1pt]{\scriptsize #2}\end{tabular}}
\SetWatermarkText{EXPECTED -- NOT OBSERVED}
\SetWatermarkFontSize{27pt}
\SetWatermarkScale{1}
\SetWatermarkColor[gray]{0.92}
\SetWatermarkAngle{45}
\pagestyle{fancy}\fancyhf{}
\fancyhead[L]{\small Know Who You Are \enspace / \enspace Expected experiment tables}
\fancyhead[R]{\small v5 \enspace 9 October 2026}
\fancyfoot[L]{\footnotesize Handwritten planning values. No benchmark results.}
\fancyfoot[R]{\small\thepage}
\setlength{\headheight}{14pt}
\setlength{\parindent}{0pt}\setlength{\parskip}{5pt}
\renewcommand{\arraystretch}{1.22}
\begin{document}
\begin{center}{\Large\bfseries Expected Experiment Tables}\end{center}
\small
Point values are conditional expectations for the specified populations. Brackets are
\textbf{subjective planning ranges}, not measured confidence intervals or calibrated
forecasts. Utility is $U=Y-0.10C$. Complete cost $C$ uses a common uniform-reference
unit; $Y$ is future-task partial quality, not native benchmark pass rate.

\begin{table}[H]
\centering
\caption*{Main Table 6 --- Expected primary results}
\setlength{\tabcolsep}{5pt}
\begin{tabularx}{\linewidth}{@{}Yccc@{}}
\toprule
Endpoint & \rare{} & Same-info. baseline & Difference / planning range\\
\midrule
Future quality & 0.608 & 0.588 & \stack{+0.020}{[-0.016, +0.056]}\\
Complete cost & 1.34 & 1.22 & \stack{+0.12}{[+0.05, +0.19]}\\
Quality--cost utility & 0.474 & 0.466 & \stack{+0.008}{[-0.035, +0.051]}\\
Assignment change & 0.20 & 0.16 & \stack{+0.04}{[-0.025, +0.105]}\\
Judgment Brier & 0.203 & 0.211 & \stack{$-0.008$}{[-0.020, +0.004]}\\
Update p95 / state bytes & 68 / 36,864 & 52 / 28,672 & \stack{+16 / +8,192}{latency: [+2, +42] ms}\\
Old-root forgetting & 0.025 & 0.041 & \stack{$-0.016$}{[-0.040, +0.008]}\\
UNKNOWN / false attribution & 0.498 / 0.0066 & 0.498 / 0.0066 & \stack{0 / 0}{identical shared-gate inputs}\\
\bottomrule
\end{tabularx}
\end{table}
\footnotesize
The baseline is contextual-trust-linear throughout, with the same responsibility gate.
Quality, cost and assignment are ArtifactRole endpoints; timing and old-root drift are
PeerSelect endpoints. Brier uses a common binary producer-contract cohort. Attribution
uses the common mutation mixture on page 2; it is not online outcome missingness.
Update p95 is the CPU feedback-processing path in milliseconds with cached features,
excluding LLM execution. State bytes are selector state, not total RAM. Assignment
change is diagnostic. Same-gate safety is an equality check, not a RARE advantage.

\small
\begin{table}[H]
\centering
\caption*{Main Table 7 --- Expected evaluation--update factorial}
\begin{tabular}{@{}llcc@{}}
\toprule
Evaluation & Update & Future utility (planning range) & Cost\\
\midrule
Masked & Off & 0.430 \rng{0.385, 0.475} & 1.02\\
Visible & Off & 0.448 \rng{0.401, 0.495} & 1.18\\
Masked & On & 0.455 \rng{0.405, 0.505} & 1.15\\
Visible & On & 0.474 \rng{0.421, 0.527} & 1.34\\
\bottomrule
\end{tabular}
\end{table}
\footnotesize
The visible/on cell is the same RARE population as the primary utility row.
The four quality means are 0.532, 0.566, 0.570 and 0.608, respectively.
Observation given updating: $0.474-0.455=0.019$; updating given observation:
$0.474-0.448=0.026$. The interaction is only
$0.474-0.448-0.455+0.430=0.001$, so these expectations do not posit a synergy gain.
The conservative interaction range implied by the four marginal ranges is
$[-0.194,0.196]$; no narrow paired interval is inferred without covariance.

\vspace{3mm}\small
The favorable-condition utility edge is small and uncertain; RARE has higher
ArtifactRole cost and higher PeerSelect latency and state. Sparse support reverses its comparison with pooled control, and
high load reverses its payoff comparison with contextual trust (page 2).
Stronger controls remain near ties: producer-contract Brier targets are 0.207 for
task-conditioned ridge and 0.206 for matched composition, versus 0.203 for RARE.

\newpage
\begin{center}{\large\bfseries Experiment matrix: expected cell results}\end{center}
\footnotesize
Row IDs match the nine-cell manuscript matrix. Each result below refers to its
specified cohort; different task or stress populations are never pooled.
\begin{table}[H]
\centering\footnotesize
\caption*{Main Table 4 --- Expected experimental matrix}
\setlength{\tabcolsep}{3pt}\renewcommand{\arraystretch}{1.30}
\begin{tabularx}{\linewidth}{@{}L{.065\linewidth}L{.13\linewidth}L{.20\linewidth}L{.19\linewidth}Y@{}}
\toprule
Cell & Track / root & Strategy / intervention & Matched comparison & Measurement / expected result\\
\midrule
RQ1-I & ArtifactRole / dev & Raw vs. terminal vs. \rare{} & Same delivery, recipient observation, legal feedback & Brier: .225 / .220 / .203. Log loss: .637 / .625 / .584. ECE: .073 / .068 / .055.\\
RQ1-II & ArtifactRole / confirm & Contextual trust vs. \rare{} & Same read cut, delay, propensity, model, budget & Brier: .211 / .203; held-out-root Brier: .231 / .217. Transfer deterioration: .020 / .014.\\
RQ2-I & ArtifactRole / confirm & No evidence vs. public evidence & Preview and assignment before target execution & Change: .04 / .12; quality: .532 / .566; adoption: .48 / .54; rework: .30 / .24.\\
RQ2-II & ArtifactRole / confirm & Delayed-only vs. evidence+delay & Same target and selected-only outcome & Utility: .455 / .474; cost: 1.15 / 1.34. Added-evidence utility: +.019.\\
RQ2-III & ArtifactRole / confirm & \rare{} vs. pooled controller & Same capacity and candidate menu & Stable utility: .474 / .472; max share: .38 / .40; min propensity: .08 / .06. Sparse utility: .427 / .450.\\
RQ3-I & PeerSelect / dev stream & Uniform, no/history, updater & Pinned graph and call/token budget & Normal payoff: .500 / .532 / .565 / .606. RARE update p50/p95: 37/68 ms; state: 36,864 B. High-load payoff RARE/context: .542/.558; p95: 480/290 ms.\\
RQ3-II & PeerSelect / drift & Update with/without replay & Fixed peer/task drift and delay & Correction / none: old-root drop .025/.048; restricted mean recovery 6.5/9.0 updates; recovery by 20 updates .96/.88; backlog p95 1.4/.9; correction/replay agreement 1.00/.86.\\
RQ4-I & ArtifactRole / eligible roots & Responsibility mutations & Producer, recipient, mixed, unknown, failure & Known-owner mixture: UNKNOWN .498; coverage .502; FA/all .0066; FA/claimed .01315; UNKNOWN precision .9592. Source-publication no-op 1.00; unknown-owner invalid credits 0.\\
RQ4-II & ArtifactRole / eligible roots & Remove one mechanism & Gate, evidence, delay, locality, updater & Utility full/gate removed/evidence removed/updater off/locality removed: .474/.461/.455/.448/.470. FA full/gate removed: .0066/.1600. Utility interaction: +.001.\\
\bottomrule
\end{tabularx}
\end{table}
\footnotesize
\textbf{Safety denominator.} The fixed owner-truth mixture weights are .50 producer-only,
.20 recipient-only, .10 mixed, .10 sink-only and .10 resource failure. Their claim coverage
is .90/.04/.18/.06/.20; all-attempt false attribution is .008/.004/.012/.002/.004.
Hence coverage is .502, UNKNOWN is .498 and false attribution is .0066; conditional
claimed-case error is $.0066/.502=.01315$. Unknown-owner cases are checked separately
and do not supply invented ground-truth labels. UNKNOWN precision uses independently required abstention, not hidden owner identity;
source-publication no-op is a parameter invariant. Nonzero FA means verdict error,
not prohibited source-credit updating. Updater-off equals the visible/off factorial
cell. Information/update ablations retain identical safety on this fixed corpus.

\textbf{Service and sampling.} Normal selected payoff is $(r+0.3)/0.6$ for the pinned
$[-0.3,0.3]$ reward contract, with dropped/timed-out arrivals scored zero in the ITT mean; it is
not mutual-cooperation rate. Recovery concerns new-task payoff (95\% of a frozen
reference for five windows, capped at 20 updates); forgetting uses old-root tasks. The online path uses cached
features, three local candidates, four logical CPU cores and a 1 GiB RAM budget. The
planning design has three roots and six reset streams per root/arm; roots are not
interchangeable independent event replicates. Workload assumptions and numeric rationale
are versioned with the JSON companion. These tables set comparison expectations, not
criteria for choosing or discarding actual experiment outcomes.
\end{document}
